\documentclass[11pt]{article}
\usepackage[margin=1in]{geometry}

% Core packages
\usepackage{amsmath,amssymb}
\usepackage[T1]{fontenc}
\usepackage[utf8]{inputenc}
\usepackage{hyperref}
\usepackage{ragged2e}

% Tighter paragraphs
\setlength{\parindent}{0pt}
\setlength{\parskip}{0.7\baselineskip}

\begin{document}

% ── Title ─────────────────────────────────────────────────────────────────
\begin{center}
{\Large\bfseries
\setlength{\baselineskip}{20pt}
Latent Heterogeneity in Health Care Utilization
Counts Using Negative Binomial Mixtures
\par}
\end{center}

% ── Authors and affiliations ──────────────────────────────────────────────
\begin{flushleft}
Katy Freitas\textsuperscript{1,2},
Susana Faria\textsuperscript{1}
and Thiago Pavin\textsuperscript{2}

\vspace{0.15cm}

{\small
\begin{tabular}{@{}ll@{}}
\textsuperscript{1} & Department of Mathematics, University of Minho, Portugal \\
\textsuperscript{2} & Independent Researcher, Health Management and Data Science, Brazil \\
\end{tabular}
}

\vspace{0.3cm}

E-mail addresses:
\href{mailto:katygdf@gmail.com}{\textit{katygdf@gmail.com}};
\href{mailto:sfaria@math.uminho.pt}{\textit{sfaria@math.uminho.pt}};
\href{mailto:thiagopavin@gmail.com}{\textit{thiagopavin@gmail.com}}

\end{flushleft}

\noindent\rule{\textwidth}{0.4pt}

% ── Abstract ──────────────────────────────────────────────────────────────
\justifying

\noindent\textbf{Abstract}

\vspace{-0.2cm}

We assess whether finite Negative Binomial mixture models recover latent
utilization groups more accurately than clustering in overdispersed count data.
We generated a synthetic dataset of 5,000 beneficiaries with five health-care utilization
variables, structured into four latent groups and calibrated to Brazilian
regulatory benchmarks. Model parameters were estimated by maximum likelihood using Expectation-Maximization algorithm.
The results suggest that finite Negative Binomial mixture models provide a more suitable approach to identifying latent utilization profiles in overdispersed health-care count data.

\noindent\textbf{Keywords:} finite mixture models, negative binomial
distribution, health-care utilization.

\noindent\rule{\textwidth}{0.4pt}

% ── Paragraph 1 ──────────────────────────────────────────────────────────
Health-care utilization data are typically characterized by discreteness,
asymmetry and overdispersion, often exhibiting substantial variability
between individuals with low and high levels of service use.
These features challenge the adequacy of purely geometric clustering
approaches, such as K-means, particularly when the underlying population
comprises latent subgroups.
Finite mixture models provide a probabilistic framework for representing
unobserved heterogeneity by assuming that the observed population arises
from a combination of latent components~\cite{mclachlan2000finite,ng2019mixture}.

% ── Paragraph 2 ──────────────────────────────────────────────────────────
In this work, we investigate whether a finite mixture model based on
Negative Binomial components can recover latent utilization profiles more
accurately than K-means.
We generated synthetic ~\cite{tucker2020generating} to emulate
health-care utilization in Brazil, following the simulation
approach of Deb and Trivedi~\cite{deb1997demand} for modelling health-care
demand via finite mixtures.
The data consists of 5,000 beneficiaries, structured into four latent
groups and described by five count variables: \textit{medical consultations}, 
\textit{emergency visits}, \textit{diagnostic exams}, \textit{hospitalizations} 
and \textit{therapies}.
The marginal distributions of the simulated variables were calibrated using
aggregate indicators from the Brazilian National Regulatory Agency for
Private Health Insurance and Plans.
Since the true latent group membership is known by construction, the
performance of each clustering method can be assessed using clustering performance was evaluated using the Adjusted Rand Index (ARI) ~\cite{hubert1985comparing} and Classification accuracy ~\cite{munkres1957algorithms}.

% ── Paragraph 3 ──────────────────────────────────────────────────────────
The proposed model assumes that each latent component follows a product
of independent Negative Binomial distributions~\cite{cameron2013regression},
one for each utilization variable.
Hospitalization and therapy variables exhibit structural excess of
zeros~\cite{lambert1992zero}, accommodated within each component.
Parameter estimation was conducted using maximum likelihood via the Expectation-Maximization
algorithm~\cite{dempster1977maximum}, with the number of components selected
according to the Bayesian Information Criterion (BIC).
For comparison, K-means clustering was applied to log-transformed and
standardized count data, with the number of clusters selected using the
Silhouette criterion. These preprocessing steps are necessary because
K-means relies on Euclidean distances, which are sensitive to scale
differences and distributional skewness \cite{hastie2009elements}.

% ── Paragraph 4 ──────────────────────────────────────────────────────────
The best performance was obtained by the Negative Binomial mixture model
when all five variables were included.
In this specification, the BIC selected four
components, matching the true number of latent groups.
The model attained ARI\,=\,0.729,
NMI\,=\,0.594, accuracy\,=\,0.877
and purity\,=\,0.861, outperforming K-means, which achieved
ARI\,=\,0.673 and accuracy\,=\,0.815.
The inclusion of rare but informative variables, particularly therapies,
improved the recovery of the latent structure.

% ── Paragraph 5 ──────────────────────────────────────────────────────────
These results provide evidence that, in overdispersed count data
with latent heterogeneity, Negative Binomial mixture models constitute a
more suitable framework than traditional clustering methods for recovering
hidden utilization profiles.
This is a methodological study and does not aim to classify real
beneficiaries or detect operational anomalies.

% ── Acknowledgements ─────────────────────────────────────────────────────
\textbf{Acknowledgements.} The research is supported
by the Auditoria Inteligente 360 project under a technological
co-development partnership agreement.

% ── References ───────────────────────────────────────────────────────────
{\small
\begin{thebibliography}{10}

\bibitem{mclachlan2000finite}
G.~J. McLachlan and D.~Peel.
\textit{Finite Mixture Models}.
John Wiley \& Sons, 2000.

\bibitem{ng2019mixture}
S.-K. Ng, L.~Xiang, and K.~K.~W. Yau.
\textit{Mixture Modelling for Medical and Health Sciences}.
Chapman and Hall/CRC, 2019.

\bibitem{tucker2020generating}
A.~Tucker, Z.~Wang, Y.~Rotalinti, and P.~Myles.
Generating high-fidelity synthetic patient data for assessing machine
learning healthcare software.
\textit{NPJ Digital Medicine}, 3(1):147, 2020.

\bibitem{deb1997demand}
P.~Deb and P.~K. Trivedi.
Demand for medical care by the elderly: a finite mixture approach.
\textit{J.\ Appl.\ Econom.}, 12(3):313--336, 1997.

\bibitem{cameron2013regression}
A.~C. Cameron and P.~K. Trivedi.
\textit{Regression Analysis of Count Data}, 2nd~ed.
Cambridge University Press, 2013.

\bibitem{lambert1992zero}
D.~Lambert.
Zero-inflated Poisson regression, with an application to defects in
manufacturing.
\textit{Technometrics}, 34(1):1--14, 1992.

\bibitem{dempster1977maximum}
A.~P. Dempster, N.~M. Laird, and D.~B. Rubin.
Maximum likelihood from incomplete data via the EM algorithm.
\textit{J.\ R.\ Stat.\ Soc.\ Ser.\ B}, 39(1):1--38, 1977.

\bibitem{hubert1985comparing}
L.~Hubert and P.~Arabie.
Comparing partitions.
\textit{J.\ Classification}, 2(1):193--218, 1985.

\bibitem{munkres1957algorithms}
J.~Munkres.
Algorithms for the assignment and transportation problems.
\textit{J.\ Soc.\ Ind.\ Appl.\ Math.}, 5(1):32--38, 1957.

\bibitem{hastie2009elements}
T.~Hastie, R.~Tibshirani, and J.~Friedman.
\textit{The Elements of Statistical Learning}, 2nd~ed.
Springer, 2009.

\end{thebibliography}
}

\end{document}